\bigskip
\noindent Autres pistes de l'onglet Findings pour ce thème, dans des cotes qui n'ont pas de section ici :\par

\begin{encadre}
\begin{theoreme*}[Grothendieck, pages 86 à 88, finitude de la normalisation\piste{23-finite-normalisation-quotient-of-regular}]\nonverifie
Soit $A$ un anneau noethérien intègre quotient d'un anneau régulier. La clôture intégrale de $A$ est finie sur $A$ si et seulement si les premiers $\mathfrak p$ de hauteur un tels que $A_{\mathfrak p}$ n'est pas régulier sont en nombre fini et que, pour chacun, la normalisée de $A_{\mathfrak p}$ est finie sur $A_{\mathfrak p}$.
\end{theoreme*}
\end{encadre}
{\small Les feuillets en esquissent une démonstration incomplète : ils construisent $C = i_*i^*\tilde B$ à travers un fermé de codimension $\geq 2$ et la montrent $(S_2)$ et $(R_1)$ par le critère de Serre ; l'universelle caténarité de $A$ vient de l'édition, et la cohérence de $C$ repose sur une phrase dont le milieu est illisible.\par}

\begin{encadre}
\begin{theoreme*}[Grothendieck, page 22, quand $\mathrm{Ind}(C)$ est un topos\piste{37-ind-topos-exercise}]\nonverifie
Pour une petite catégorie $C$ d'enveloppe de Karoubi $C'$, $\mathrm{Ind}(C)$ est un topos si et seulement si $C'$ a les limites inductives finies et que son foncteur canonique vers les faisceaux pour la topologie finitaire les respecte ; ce topos est alors parfait exactement quand $C'$ a les limites projectives finies.
\end{theoreme*}
\end{encadre}
{\small Mathlib n'a pas les conditions de SGA~4 pour que $\mathrm{Ind}(C)$ soit un topos. On se reportera aux notes de Grothendieck (\lot{37}{2}{22}).\par}

\begin{encadre}
\begin{conjecture*}[Grothendieck, pages 1 à 4, Dold--Kan hors de $\Delta$\piste{112-linearization-beyond-the-simplex}]
Pour une petite catégorie $A$, la restriction aux préfaisceaux abéliens \q{lisses}, ceux le long desquels la linéarisation ne varie pas, donne une correspondance de type Dold--Kan avec une catégorie dérivée de complexes localement constants. La page pose la question sous plusieurs formes et relève que $L_b$ porte des coefficients tordus et non constants.
\end{conjecture*}
\end{encadre}
{\small Mathlib n'a pas les catégories dérivées de complexes localement constants. On se reportera aux notes de Grothendieck (\lot{112}{1}{1}).\par}

\begin{encadre}
\begin{theoreme*}[Grothendieck, pages 6 et 7, recollement de topos de préfaisceaux\piste{116-presheaf-gluing-continuous-gluing-functor}]\nonverifie
Un topos recollé d'un sous-topos ouvert $X_0$ et de son fermé complémentaire $X_1$ le long d'un foncteur exact à gauche $\varphi$ est un topos de préfaisceaux exactement quand $X_0$ et $X_1$ sont de la forme $\mathrm{Top}(\mathcal X_0)$ et $\mathrm{Top}(\mathcal X_1)$ et que $\varphi$ commute aux limites projectives quelconques ; la donnée de recollement est alors un bifoncteur $\mathcal X_0^{\mathrm{op}} \times \mathcal X_1 \to (\mathrm{Ens})$.
\end{theoreme*}
\end{encadre}
{\small Mathlib n'a pas le recollement des topos. On se reportera aux notes de Grothendieck (\lot{116}{1}{6}).\par}
